\subsubsection{口交、乳交、肛交的安全细节补遗}

本书与女性卷都已系统介绍过这三种方式（参见《口交的艺术与亲密意义》章、《口交技巧》与《肛交技巧与体验》等章节）。这里只补三处前文着墨较少的细节：\textbf{口交的感染风险清单、肛交的准备与疼痛处理、乳交的实操与哺乳期边界}。

\textbf{一、口交：被低估的感染清单}
"口交比性交安全"是流传最广的半真半假的说法。它的确不涉及怀孕风险，但传播感染的能力并不低——关键在于\textbf{双方黏膜的破损情况}。
\begin{itemize}
\item \textbf{咽部淋病}：口腔与咽喉的淋球菌感染常常无症状，却具有传染性；表现为咽痛时，容易被当作普通咽炎自行用药，导致漏诊与耐药；
\item \textbf{口腔与生殖器 HPV}：口交与口腔 HPV 感染相关，而口腔 HPV 又与口咽部肿瘤相关（风险显著低于吸烟与饮酒，但确实存在）；HPV 疫苗对已开始的性活动仍有部分保护价值；
\item \textbf{疱疹（HSV-1/2）}：口唇疱疹可在口交中传播到生殖器，反之亦然——这是生殖器疱疹常见的感染途径之一，也是"我从来没有过不洁性行为却感染了"的常见解释；
\item \textbf{其他}：梅毒（口腔硬下疳易被忽略）、衣原体、乙肝（精液与阴道分泌物可传播）。HIV 经口交传播的效率低于肛交与阴道性交，但在口腔破损、牙龈出血或存在其他性传播感染时风险上升；
\item \textbf{可执行做法}：口交时使用避孕套（男性）或口腔屏障膜（dental dam，可剪开避孕套替代）；避免在口腔溃疡、牙龈出血、近期拔牙或咽部手术期间进行；把口腔检查纳入常规体检（尤其长期吸烟饮酒者）。\textbf{吞咽精液本身不是感染途径的关键}——风险来自黏膜接触，而非"咽下去"这个动作。
\end{itemize}

\textbf{二、肛交：准备、疼痛与界限}
\begin{itemize}
\item \textbf{准备}：肠道本身不需要"深度清洁"——频繁灌肠反而会破坏直肠黏膜的保护层、增加损伤与感染风险。性行为前正常排便、清洁外肛周即可；如确有需要，浅度清洁（不超过外括约肌以上）优于深灌肠，且不建议频繁进行；
\item \textbf{润滑是必需品而非加成}：肛门与直肠不分泌润滑液，必须使用足量\textbf{水性或硅基润滑剂}；油性润滑剂（凡士林、婴儿油）会腐蚀乳胶安全套，应避免；
\item \textbf{疼痛的处理}：\textbf{肛交不应以"疼痛是正常的"为前提}。疼痛通常是三个信号：润滑不足、进入过快、或外括约肌尚未放松。正确做法是退回一步、补充润滑、延长停顿，而不是"忍一下就过去了"；持续或剧烈疼痛应考虑肛裂（多与机械损伤相关）并就医；
\item \textbf{顺序规则}：\textbf{肛交后的器具或阴茎不得直接进入阴道或口腔}——这一条能预防大量感染性事件（包括细菌性阴道病与肠道菌移位引起的尿路感染）；
\item \textbf{屏障}：肛交是 HIV 传播效率最高的性行为之一，全程安全套是底线；口腔屏障膜、指套同样适用于其他接触方式；
\item \textbf{禁忌与就医信号}：有肛裂、痔急性发作、炎症性肠病活动期、近期肛门手术者应避免或先咨询医生；出现出血、持续疼痛、发热或分泌物时及时就诊。
\end{itemize}

\textbf{三、乳交：实操与哺乳期边界}
\begin{itemize}
\item \textbf{实操要点}：乳交（乳房夹持刺激）主要依赖乳沟的包裹感与乳房柔软度，个人体型差异很大，不必追求"标准做法"；充分润滑可减少摩擦不适；与口交结合是常见变体（前文已有介绍）；
\item \textbf{与哺乳期的边界}：哺乳期乳房可能出现泌乳（受刺激时溢乳属正常生理反应），但这不意味着可以从"非婴儿的成年人口中吮吸"中获得额外收益——\textbf{成人在哺乳期吸吮不会增加泌乳量，也不是催乳疗法}；同时要注意哺乳期乳房的胀痛与感染风险（乳腺炎）；
\item \textbf{疼痛与肿块的红旗}：任何单侧的、固定位置的乳房肿块，或持续的乳头溢液（尤其是血性或单侧），都应尽早就医检查——这一点同样适用于男性；
\item \textbf{关于乳房刺激的沟通}：乳房与乳头是敏感区域，但敏感度与偏好个体差异极大；对许多人而言，"被碰触的方式"比"碰触的部位"更关键。这仍是本书反复强调的那句话：\emph{问到具体的偏好，比猜测更有用。}
\end{itemize}
